\documentclass{scrreprt}

\usepackage{pylab}

\begin{document}
	
	\myheader{Методичні вказівки}
	
	{\large
		\section*{Лабораторна робота: Умовні оператори в Python}
		
		При виконанні лабораторних робіт з теми \textbf{<<Умовні оператори в Python>>} зверніть увагу на такі рекомендації:
		
		\begin{enumerate}[label=\arabic*.]
			\item \textbf{Типи даних та введення:}
			\begin{itemize}
				\item Для натуральних та цілих чисел використовуйте \texttt{int(input())}.
				\item Для дійсних чисел та координат використовуйте \texttt{float(input())}.
			\end{itemize}
			\item \textbf{Форматування виведення:}
			\begin{itemize}
				\item Якщо в умові вимагається точність до сотих, використовуйте f-рядки: \texttt{f"\{val:.2f\}"} або функцію \texttt{round(val, 2)}.
			\end{itemize}
			\item \textbf{Логічні оператори:}
			\begin{itemize}
				\item Застосовуйте \texttt{and} (логічне І), \texttt{or} (логічне АБО), \texttt{not} (логічне НІ).
			\end{itemize}
			\item \textbf{Обробка вироджених та граничних випадків:}
			\begin{itemize}
				\item У геометрії завжди перевіряйте збіг точок (наприклад, $A = B$), ділення на нуль, існування трикутника (нерівність трикутника $a + b > c$) та колінеарність точок.
			\end{itemize}
		\end{enumerate}
	}
	
	\newpage
	
	\myheader{Варіант 1}
	
	{\large
		\begin{enumerate}
			\item \textbf{Числа, що закінчуються на 1.} Дано два натуральних числа. Надрукувати <<Так>>, якщо обидва закінчуються цифрою 1. В іншому разі надрукувати <<Ні>>.
			\item \textbf{Цеглина у круглому отворі.} Є цеглина у формі прямокутного паралелепіпеда з лінійними розмірами $a$, $b$ та $c$~см і круглий отвір у стіні діаметром $d$~см. Чи пройде цеглина в цей отвір?
			\item \textbf{Дослідження СЛАР $2\times 2$.} Дано довільні числа $a_1, b_1, c_1, a_2, b_2, c_2$. Дослідити розв'язність системи рівнянь
			$$
			\begin{cases}
				a_1 x + b_1 y = c_1, \\
				a_2 x + b_2 y = c_2.
			\end{cases}
			$$
			Відповіддю до задачі мають бути повідомлення: <<Система несумісна>>, <<Система має безліч розв'язків>> або <<Система має єдиний розв'язок>>. У останньому випадку знайти цей розв'язок із точністю до сотих.
		\end{enumerate}
	}
	
	\vspace{2cm}
	
		
	\myheader{Варіант 2}
	
	{\large
		\begin{enumerate}
			\item \textbf{Перестановка цифр трицифрових чисел.} Дано два трицифрових натуральних числа. Чи можна з цифр одного числа скласти інше? Розглянути всі можливі випадки.
			\item \textbf{$k$-та цифра послідовності $10\dots 99$.} Нехай у послідовності \texttt{10111213...979899} виписано поспіль усі двоцифрові натуральні числа. Користувач вводить натуральне число $k$ ($1 \leqslant k \leqslant 180$). Надрукувати $k$-ту цифру цієї послідовності.
			\item \textbf{Дослідження рівнобедреного трикутника.} Точки $A, B, C$ на площині задано своїми координатами. Визначити, чи є трикутник $ABC$ рівнобедреним. Розглянути всі можливі випадки розташування точок (збіг будь-яких двох точок; усі точки лежать на одній прямій). У разі виродженого трикутника виводити відповідне повідомлення.
		\end{enumerate}
	}
	
	\newpage
	
	\myheader{Варіант 3}
	
	{\large
		\begin{enumerate}
			\item \textbf{Не менше двох непарних чисел.} Дано три натуральних числа. Надрукувати <<Так>>, якщо хоча б два з них непарні. В іншому разі надрукувати <<Ні>>.
			\item \textbf{Моделювання роботи світлофора.} Робота світлофора для водіїв запрограмована таким чином: на початку кожної години протягом хвилини горить зелений сигнал, потім протягом трьох секунд "--- жовтий, протягом 27 секунд "--- червоний, протягом хвилини "--- знову зелений і т.\,д. Дано число $t$, що означає час у секундах, який минув від початку чергової години. Визначити, сигнал якого кольору горить для водіїв у цей момент.
			\item \textbf{Належність двох точок прямій.} Дано 6 довільних чисел: координати точок $(x_1, y_1)$ та $(x_2, y_2)$ на площині, а також коефіцієнти $k$ і $b$ рівняння прямої $y = kx + b$. Чи вірно, що обидві точки лежать на даній прямій і не збігаються?
		\end{enumerate}
	}
	
	\vspace{2cm}
	
	\myheader{Варіант 4}
	
	{\large
		\begin{enumerate}
			\item \textbf{Мінімальне число є непарним.} Дано три натуральних числа. Надрукувати <<Так>>, якщо мінімальне з них непарне. В іншому разі надрукувати <<Ні>>.
			\item \textbf{Обчислення плати за телефон.} Послуги телефонної мережі оплачуються за таким правилом: за розмови до $A$ хвилин на місяць включно сплачується фіксована сума $B$~грн, а розмови понад установу норму оплачуються за тарифом $C$~грн за хвилину. Окрім тарифів, дано ціле число "--- кількість хвилин розмов за минулий місяць. Обчислити плату за користування телефоном.
			\item \textbf{Взаємне розташування двох прямих.} Дано чотири довільних числа $k_1, b_1, k_2, b_2$. Визначити взаємне розташування прямих $y = k_1 x + b_1$ та $y = k_2 x + b_2$ (паралельні, перетинаються або збігаються). Урахувати всі можливі випадки.
		\end{enumerate}
	}
	
	\newpage
	
	\myheader{Варіант 5}
	
	{\large
		\begin{enumerate}
			\item \textbf{Наявність від'ємного числа.} Дано три цілих числа. Надрукувати <<Так>>, якщо хоча б одне з них від'ємне. В іншому разі надрукувати <<Ні>>.
			\item \textbf{Колінеарність трьох точок.} Дано координати точок $A, B, C$ на площині. Визначити, чи лежать ці точки на одній прямій. Урахувати всі можливі випадки розташування точок.
			\item \textbf{Паралельність прямої осям координат.} Дано координати точок $A$ та $B$ на площині. Визначити, чи паралельна пряма $AB$ одній із координатних осей, і якщо так, то якій саме. У разі збігу точок вивести повідомлення <<Пряму побудувати неможливо>>.
		\end{enumerate}
	}
	
	\vspace{2cm}
	
	\myheader{Варіант 6}
	
	{\large
		\begin{enumerate}
			\item \textbf{Кратність 4 без 0 на кінці.} Дано натуральне число. Надрукувати <<Так>>, якщо це число кратне 4 і не закінчується нулем. В іншому разі надрукувати <<Ні>>.
			\item \textbf{Периметр і площа трикутника.} Дано три довільних числа $a, b, c$. Якщо вони можуть бути довжинами сторін трикутника, знайти його периметр і площу з точністю до сотих. В іншому разі надрукувати повідомлення про помилку.
			\item \textbf{Взаємне розташування прямої та кола.} Дано довільні числа $a, b, c, x_0, y_0, r$. Дослідити взаємне розташування прямої $ax + by + c = 0$ та кола $(x - x_0)^2 + (y - y_0)^2 = r^2$:
			\begin{enumerate}[label=\alph*)]
				\item дотикаються (одна спільна точка);
				\item перетинаються (дві спільні точки);
				\item не мають спільних точок.
			\end{enumerate}
		\end{enumerate}
	}
	
	\newpage
	
	\myheader{Варіант 7}
	
	{\large
		\begin{enumerate}
			\item \textbf{Кратність добутку сумі двох чисел.} Дано три цілих числа. Надрукувати <<Так>>, якщо добуток усіх трьох чисел кратний сумі будь-яких двох із них. В іншому разі надрукувати <<Ні>>.
			\item \textbf{Обчислення повного віку людини.} Відомі рік, номер місяця та день народження людини, а також рік, номер місяця та номер поточного дня місяця. Визначити вік людини (кількість повних років).
			\item \textbf{Включення одного круга в інший.} Визначити, чи лежить один із кругів, заданих рівняннями $(x - x_1)^2 + (y - y_1)^2 \leqslant r_1^2$ та $(x - x_2)^2 + (y - y_2)^2 \leqslant r_2^2$, повністю всередині іншого. Числа $x_i, y_i, r_i$ ($i = 1, 2$) вводяться з клавіатури. Випадок збігу кругів не розглядати.
		\end{enumerate}
	}
	
	
	\vspace{2cm}
	\myheader{Варіант 8}
	
	{\large
		\begin{enumerate}
			\item \textbf{Пара чисел із найменшою відстанню.} Дано три цілих числа. Вивести ті два з них, відстань між якими на числовій осі є найменшою.
			\item \textbf{Перевірка перебування потяга.} Потяг прибуває на станцію о $h_1$ годин $m_1$ хвилин, а відправляється о $h_2$ годин $m_2$ хвилин. Пасажир прийшов на платформу о $x$ годин $y$ хвилин. Чи стоятиме потяг на платформі? Усі числа вводяться з клавіатури. Під час введення перевіряти дані на коректність.
			\item \textbf{Належність початку координат трикутнику.} Дано координати точок $A, B, C$ на площині. Визначити, чи лежить початок координат усередині трикутника $ABC$.
		\end{enumerate}
	}
	
	\newpage
	
	\myheader{Варіант 9}
	
	{\large
		\begin{enumerate}
			\item \textbf{Наявність числа, кратного 5.} Дано три натуральних числа. Надрукувати <<Так>>, якщо хоча б одне з них кратне 5. В іншому разі надрукувати <<Ні>>.
			\item \textbf{Касовий апарат: розрахунок решти.} Скласти програму, яка запитує в користувача вартість покупки та суму грошей, унесену покупцем.
			\begin{itemize}
				\item Якщо решта не потрібна, надрукувати на екрані <<Дякуємо>>;
				\item Якщо внесено більше грошей, ніж потрібно, надрукувати <<Візьміть решту ... грн>> (із зазначенням потрібної суми);
				\item Якщо ж грошей недостатньо, надрукувати <<Для покупки внесіть ще ... грн>> також із зазначенням необхідної суми.
			\end{itemize}
			\item \textbf{Перевірка точки як середини відрізка.} Точки $A, B, C$ на площині задано своїми координатами. Визначити, чи є точка $C$ серединою відрізка $AB$.
		\end{enumerate}
	}
	
	\vspace{2cm}
	
	\myheader{Варіант 10}
	
	{\large
		\begin{enumerate}
			\item \textbf{Геометрична прогресія з цифр.} Дано трицифрове натуральне число. Визначити, чи можна за якоїсь перестановки цифр цього числа (включно з тотожною перестановкою) отримати геометричну прогресію.
			\item \textbf{Визначення дня тижня для $k$-го дня.} Дано натуральне число $k$ ($1 \leqslant k \leqslant 365$). Визначити, яким днем тижня буде $k$-й день року, якщо 1~січня цього року "--- п'ятниця.
			\item \textbf{Дослідження системи з модулем $|x|+|y|=1$.} Дано довільні числа $a$ та $b$. Дослідити розв'язність системи рівнянь
			$$
			\begin{cases}
				y = ax + b, \\
				|x| + |y| = 1.
			\end{cases}
			$$
			У відповіді вивести одне з повідомлень: <<Система не має розв'язків>>, <<Система має єдиний розв'язок>>, <<Система має два розв'язки>>, <<Система має безліч розв'язків>>.
		\end{enumerate}
	}
	
	\newpage
	
	\myheader{Варіант 11}
	
	{\large
		\begin{enumerate}
			\item \textbf{Відсутність чисел із 0 на кінці.} Дано три натуральних числа. Надрукувати <<Так>>, якщо серед них немає числа, що закінчується нулем. В іншому разі надрукувати <<Ні>>.
			\item \textbf{Класифікація трикутника за сторонами.} Дано три довільних числа $a, b, c$. Чи можуть вони бути довжинами сторін деякого трикутника? Якщо так, визначити його тип (рівносторонній, рівнобедрений, прямокутний, загального вигляду).
			\item \textbf{Перебування кола в координатній чверті.} Визначити, чи лежить коло, задане рівнянням $(x - a)^2 + (y - b)^2 = r^2$, повністю в якійсь із координатних чвертей. Числа $a, b, r$ вводяться з клавіатури.
		\end{enumerate}
	}
	
	
	\vspace{2cm}
	\myheader{Варіант 12}
	
	{\large
		\begin{enumerate}
			\item \textbf{Рівно два двоцифрових числа.} Дано три натуральних числа. Надрукувати <<Так>>, якщо рівно два з них двоцифрові. В іншому разі надрукувати <<Ні>>.
			\item \textbf{$k$-та цифра послідовності $0\dots 20$.} Нехай у послідовності \texttt{0123...181920} виписано поспіль усі цілі числа від 0 до 20. Користувач вводить натуральне число $k$ ($1 \leqslant k \leqslant 32$). Надрукувати $k$-ту цифру цієї послідовності.
			\item \textbf{Перетин двох відрізків $AB$ та $CD$.} Дано координати точок $A, B, C, D$ на площині. Визначити, чи перетинаються відрізки $AB$ та $CD$. Урахувати всі можливі випадки розташування точок.
		\end{enumerate}
	}
	
	\newpage
	
	\myheader{Варіант 13}
	
	{\large
		\begin{enumerate}
			\item \textbf{Наявність хоча б двох однакових.} Дано три натуральних числа. Надрукувати <<Так>>, якщо серед них є хоча б два однакових. В іншому разі надрукувати <<Ні>>.
			\item \textbf{Порівняння віку двох людей.} Відомі дати народження (день, місяць, рік) двох людей. Передбачається, що дати не збігаються. Визначити, хто з них старший.
			\item \textbf{Обчислення кутів та типу трикутника.} Дано три довільних числа $a, b, c$. Чи можуть вони бути довжинами сторін деякого трикутника? Якщо так, знайти його кути та визначити тип (гострокутний, тупокутний, прямокутний).
		\end{enumerate}
	}
	
	\vspace{2cm}
	\myheader{Варіант 14}
	
	{\large
		\begin{enumerate}
			\item \textbf{Усі три числа є різними.} Дано три натуральних числа. Надрукувати <<Так>>, якщо серед них немає однакових. В іншому разі надрукувати <<Ні>>.
			\item \textbf{Корені кубічного рівняння $ax^3+bx=0$.} Дано два довільних числа $a$ та $b$. Знайти всі дійсні корені рівняння $ax^3 + bx = 0$. Якщо рівняння не має дійсних коренів або їх безліч, повідомити про це.
			\item \textbf{Дотик двох кіл (зовнішній/внутрішній).} Визначити, чи дотикаються кола, задані рівняннями $(x - x_1)^2 + (y - y_1)^2 = r_1^2$ та $(x - x_2)^2 + (y - y_2)^2 = r_2^2$ (дотик може бути внутрішнім або зовнішнім). Числа $x_i, y_i, r_i$ ($i = 1, 2$) вводяться з клавіатури. Випадок збігу кіл не розглядати.
		\end{enumerate}
	}
	

	\newpage
	\myheader{Варіант 15}
	
	{\large
		\begin{enumerate}
			\item \textbf{Рівно одне парне число.} Дано три натуральних числа. Надрукувати <<Так>>, якщо рівно одне з них парне. В іншому разі надрукувати <<Ні>>.
			\item \textbf{Табло аеропорту: наступна секунда.} Електронне табло аеропорту показує час $h:m:s$, де $0 \leqslant h \leqslant 23$ "--- години, $0 \leqslant m \leqslant 59$ "--- хвилини, $0 \leqslant s \leqslant 59$ "--- секунди. Який час покаже табло в наступну секунду? Вважається, що в цей момент не здійснюватиметься перехід на зимовий або літній час.
			\item \textbf{Проходження прямої через початок координат.} Точки $A$ та $B$ на площині задано своїми координатами. Визначити, чи проходить пряма $AB$ через початок координат. Якщо точки збігаються, вивести повідомлення <<Пряму побудувати неможливо>>. Також урахувати випадки розташування точок на координатних осях.
		\end{enumerate}
	}
	

\end{document}
